\documentclass[11pt]{article}
\usepackage{amsmath,amssymb,amsthm}
\usepackage[margin=1in]{geometry}
\usepackage{enumitem}

\newcommand{\E}{\mathbb{E}}
\newcommand{\1}{\mathbb{1}}
\newcommand{\Valid}{\mathrm{Valid}}

\title{Math Formulation: Constraining a Hacked Reward ---\\
Gates, Aggregation, and Why In-Loop Alerts Fail}
\author{S Hariharan, J Shrenik Reddy, Ansa N}
\date{}

\begin{document}
\maketitle

\noindent This continues from \texttt{open\_items.pdf}, which formalized
the three items left open by the working draft. Section numbers continue
from 10; equation numbers continue from (16). Where that document
specified the \emph{algorithm}, this one specifies the \emph{constraints}
added after the algorithm was found to be exploitable, and states
precisely what each does and does not guarantee. Every claim marked
``measured'' is reproduced in \texttt{documentation/reward\_hacking\_%
investigation.md}.

\section{The validity predicate as the place to spend a hard constraint}
\label{sec:gate}

Section 4 defined $\Valid(m)$ as RDKit sanitization success, and Eq.~5
collapses invalid molecules to a flat reward $r_{\mathrm{inv}}$. Write the
sanitization predicate as $\Valid_0(m)$ and generalize:
\begin{equation}
\Valid(m) \;=\; \Valid_0(m) \,\wedge\, \1\!\big[\,\mathrm{HAC}(m) \ge h_{\min}\,\big],
\label{eq:gate}
\end{equation}
with $\mathrm{HAC}$ the heavy-atom count. Everything downstream is
unchanged: gated molecules take $r_{\mathrm{inv}}$ and the hurdle
advantage floor of Section 9, exactly as a valence violation does.

\paragraph{Why a gate and not a term.} Two arguments, and they are
different. First, \emph{monotonicity}: any reward term increasing in
molecular size has its optimum at unbounded size, so a soft size reward is
hackable in the opposite direction. A hinge $\max(0, h_{\min}-\mathrm{HAC})$
avoids that, but---second---\emph{gradient conditioning}: a term intended
to drive the reward to zero must approach zero, and under the
multiplicative aggregation of Section~\ref{sec:geom} the sensitivity
$\partial G/\partial S_k \propto G/S_k$ diverges there. Eq.~\eqref{eq:gate}
has no gradient to climb and none to explode. The constraint is spent
where it is free.

\paragraph{The cost, and the condition under which it is not free.}
Gating removes candidates from the group. Writing $n_k$ for the number of
valid molecules at step $k$, Section 9 established that the group-relative
advantage carries no information when $n_k < 2$, and returns exactly zero.
So Eq.~\eqref{eq:gate} is safe only while
\begin{equation}
\Pr\big[\,n_k \ge 2\,\big] \approx 1 ,
\label{eq:starvation}
\end{equation}
which must be monitored rather than assumed. Measured at $h_{\min}=10$ and
$N=64$: $n_k \in [61,62]$ on average, and the gated fraction falls from
$\approx 9\%$ at initialization to $2$--$4\%$ by step 150. That decline is
the diagnostic signature of a hard gate: the policy learns to avoid it
rather than paying it repeatedly, which a soft penalty would not produce.

\section{Multiplicative aggregation}
\label{sec:geom}

Eq.~5 aggregates the reward terms linearly. A weighted sum permits
compensation: a term at its maximum can outbid another near its minimum,
which is precisely how a molecule scoring $C=0.999$ and $\mathrm{QED}=0.47$
outranks one scoring $C=0.811$ and $\mathrm{QED}=0.54$. Replace the sum
with a weighted geometric mean over terms $S_k \in [0,1]$:
\begin{equation}
G(m) \;=\; \Big(\textstyle\prod_k S_k(m)^{\,w_k}\Big)^{1/\sum_k w_k}
\;=\; \exp\!\Big(\tfrac{1}{\sum_k w_k}\textstyle\sum_k w_k \log S_k(m)\Big).
\label{eq:geom}
\end{equation}
The right-hand form is what should be implemented: it is stable, and it
makes the failure mode explicit, since $\log S_k \to -\infty$ as
$S_k \to 0$. Clipping $S_k$ to $[\varepsilon, 1]$ bounds both the value and
the sensitivity.

Two properties follow, and only the first is usually stated.
\begin{enumerate}[nosep]
  \item \textbf{Veto.} $G \to 0$ if any $S_k \to 0$, regardless of the
  others. This is the intended behaviour.
  \item \textbf{Dilution.} The exponent $w_k / \sum_j w_j$ attenuates every
  factor. A term contributing $S_k$ changes $G$ by $S_k^{w_k/\sum w}$,
  so with five equally weighted terms a factor of $0.55$ becomes $0.89$.
  A penalty calibrated on $S_k$ in isolation will therefore be far weaker
  in aggregate than intended, and must be calibrated on its exponentiated
  contribution instead.
\end{enumerate}
Property 2 is why Section~\ref{sec:gate} is a gate rather than a factor:
a veto expressed multiplicatively is diluted, a veto expressed through
$\Valid$ is not.

\paragraph{Scale compatibility.} Eq.~\eqref{eq:geom} requires
$S_k \in [0,1]$. Any monotone recalibration of a term that leaves that
interval---such as a logit transform centred on a reference mean, which is
negative below it---is incompatible, and the incompatibility is structural
rather than data-dependent: whether a negative value appears in a given
group is a property of the sample, not of the configuration.

\section{Why an in-loop substructure penalty does not generalize}
\label{sec:goodhart}

Let $\mathcal{A}$ be a finite set of substructure patterns and
\begin{equation}
N_{\mathcal{A}}(m) \;=\; \sum_{a \in \mathcal{A}} \1\big[\,a \sqsubseteq m\,\big],
\qquad
S_{\mathcal{A}}(m) \;=\; e^{-\lambda N_{\mathcal{A}}(m)},
\label{eq:alert}
\end{equation}
with $a \sqsubseteq m$ denoting a substructure match. Adding
$S_{\mathcal{A}}$ to Eq.~\eqref{eq:geom} penalizes members of $\mathcal{A}$
and \emph{nothing else}.

This is the crux. The quantity of interest is chemical reactivity, call it
$R^\star(m)$, which is unobserved. $N_{\mathcal{A}}$ is a proxy, and the
policy optimizes the proxy. Because $\mathcal{A}$ is finite and the space
of reactive motifs is not, there exist molecules with
$N_{\mathcal{A}}(m)=0$ and $R^\star(m)$ large, and the optimizer has a
gradient toward exactly those. Formally, optimizing $S_{\mathcal{A}}$
reduces $\E[N_{\mathcal{A}}]$ under $\pi_\theta$ with no implication for
$\E[R^\star]$ --- the two are related only through whatever correlation
$\mathcal{A}$ happens to carry, and the optimization is free to destroy it.

\paragraph{The held-out instrument criterion.} Since $R^\star$ is
unobservable, the only available test is a second pattern set
$\mathcal{B}$ with $\mathcal{B} \cap \mathcal{A} = \emptyset$, never used
in training. An intervention is credited only if
\begin{equation}
\E_{\pi_\theta}\big[N_{\mathcal{B}}\big] \;<\;
\E_{\pi_{\theta_{\mathrm{base}}}}\big[N_{\mathcal{B}}\big]
\label{eq:heldout}
\end{equation}
by a margin exceeding the seed-to-seed variation of $N_{\mathcal{B}}$.
Measured at $\lambda=0.3$, $w_{\mathcal{A}}=3$: the trained-against rate
fell $0.339 \to 0.203$ while BRENK moved $0.438 \to 0.406$ and
BRENK$+$PAINS$+$NIH $0.464 \to 0.427$, both within seed noise. Eq.~%
\eqref{eq:heldout} is not satisfied.

\paragraph{Corollary.} Training against $\mathcal{A} \cup \mathcal{B}$ to
close the gap is self-defeating: it leaves no set disjoint from the
training signal, so Eq.~\eqref{eq:heldout} becomes uncomputable and the
failure becomes undetectable rather than absent. At least one instrument
must be permanently withheld.

\section{Post-hoc selection as a distinct operator}
\label{sec:postfilter}

The argument of Section~\ref{sec:goodhart} turns on $\mathcal{A}$ entering
the objective. It therefore does not apply to a filter applied after
training. Let $\pi_\theta$ be the trained policy and define the selection
operator
\begin{equation}
\mathcal{F}\big[\pi_\theta\big](m) \;=\;
\frac{\pi_\theta(m)\,\1\big[\,N_{\mathcal{A}\cup\mathcal{B}}(m) = 0\,\big]}
     {\Pr_{\pi_\theta}\!\big[N_{\mathcal{A}\cup\mathcal{B}} = 0\big]},
\label{eq:postfilter}
\end{equation}
the policy conditioned on passing the union of every available instrument.
Since $\theta$ is fixed when $\mathcal{F}$ is applied, no gradient and no
reward reaches $\mathcal{A} \cup \mathcal{B}$, and the union may be used
in full --- the constraint of Section~\ref{sec:goodhart} is lifted exactly
because nothing is being trained.

What Eq.~\eqref{eq:postfilter} buys and what it does not:
\begin{enumerate}[nosep]
  \item It is not gameable, by construction.
  \item It does not change $\pi_\theta$. The generator remains as reactive
  as before; the denominator is the yield, measured at $0.40$.
  \item It is not a safety guarantee. $\Pr[N=0]$ is bounded below $1$ even
  for approved drugs --- measured, $39\%$ of the BBB$+$ molecules in BBBP
  fail the same union --- so the predicate is neither necessary nor
  sufficient for safety, only for the absence of \emph{recognized} motifs.
\end{enumerate}

\section{A stopping criterion from the reward's own uncertainty}
\label{sec:snr}

Section 10.3 tracked $\mathrm{Var}_i[D_{\phi_k}(m_i)]$ as a diagnostic for
the discriminator. The frozen classifier admits a sharper statement,
because its uncertainty can be estimated directly.

Let $\ell(m) = \operatorname{logit} C(m)$ and let $\tilde{C}^{(1)},\dots,
\tilde{C}^{(K)}$ be $K$ stochastic forward passes with dropout active and
normalization layers held in inference mode --- the latter being required,
or $C$ ceases to be a function of $m$ alone and the stationarity argument
of Section 3.1 fails. Define
\begin{equation}
\sigma_{\mathrm{sig}}(k)^2 = \mathrm{Var}_{i \in \mathcal{G}_k}\big[\bar\ell(m_i)\big],
\qquad
\sigma_{\mathrm{noise}}(k)^2 = \E_{i \in \mathcal{G}_k}\big[\mathrm{Var}_j\,\tilde\ell^{(j)}(m_i)\big],
\label{eq:snrdef}
\end{equation}
the between-molecule spread of the consensus logit and the mean
within-molecule spread across dropout samples, over group $\mathcal{G}_k$.
The first is what the group-relative advantage ranks on; the second is the
model's own epistemic uncertainty about each ranking. Their ratio
\begin{equation}
\mathrm{SNR}(k) \;=\; \sigma_{\mathrm{sig}}(k) \,/\, \sigma_{\mathrm{noise}}(k)
\label{eq:snr}
\end{equation}
is therefore the quantity that decides whether the reward is still
informative. When $\mathrm{SNR}(k) < 1$ the differences the advantage ranks
on are smaller than the classifier's uncertainty about them, and further
optimization against $C$ pursues noise.

Measured on the baseline: $\mathrm{SNR}$ falls $2.52 \to 1.24 \to 0.42$ at
steps $0$, $100$, $199$, crossing $1$ between steps 100 and 150 --- the
same interval over which scaffold diversity collapses.

\paragraph{Two consequences.} First, $\sigma_{\mathrm{noise}}$ is
approximately constant in $k$ (measured: $0.90$--$0.98$ throughout) while
$\sigma_{\mathrm{sig}}$ falls by $5.5\times$, so Eq.~\eqref{eq:snr} is
driven almost entirely by signal loss rather than by the model becoming
less certain. Second, the apparent saturation of $C$ in probability space
is an artifact of the sigmoid: the same MC-dropout spread is flat on the
logit scale ($0.86$ on real molecules, $0.99$ on collapsed samples) while
appearing to shrink $3.6\times$ in probability. A recalibration therefore
restores $\sigma_{\mathrm{sig}}$ but leaves $\sigma_{\mathrm{noise}}$
untouched --- it makes the reward legible, not more informative, and
Eq.~\eqref{eq:snr} is the statement of that distinction.

\section{What none of this addresses}
\label{sec:residual}

Each constraint above narrows the policy's access to a particular
exploit. None changes the objective. Writing $y(m)$ for the BBBP label,
the classifier estimates
\begin{equation}
C(m) \;\approx\; \Pr\big[\,y(m) = 1 \,\big]
\;=\; \Pr\big[\,m \text{ crosses the blood--brain barrier}\,\big],
\end{equation}
which is a statement about transport and carries no term for therapeutic
intent. The training set is consistent with that reading: of the 65
molecules with $\mathrm{HAC} \le 8$ in BBBP, 62 are labelled positive, and
they include chloroform, dichloromethane, bromoform and nitrous oxide.
Carbon tetrachloride is therefore not an extrapolation failure of $C$ but
an interpolation between positive training examples, and $C$ is correct
about the only question it was asked.

Any constraint added to Eq.~5 is a counterweight to that objective rather
than a correction of it. Five have now been tested --- property-distribution
distance, applicability domain, Tox21 activity, in-loop substructure
alerts, and multiplicative aggregation --- and only the size gate and
post-hoc selection survive their own validation. The remaining
intervention that acts on the objective itself is curation of the label
set, and it has not been attempted.

\end{document}
